
\begin{enumerate}
    \item \textbf{Descarga y organiza el conjunto de imágenes:} Puedes utilizar el dataset clásico
          en análisis de texturas Brodatz Texture Dataset, el conjunto Flickr Material Database (FMD) o
          un conjunto personalizado de distintas clases:
          \begin{itemize}
              \item corteza, ladrillo, burbujas, hierba, cuero, piel, rafia, arena, paja, agua, tejido,
                    madera, lana
          \end{itemize}
    \item \textbf{Carga las imágenes y conviértelas a escala de grises:} El Brodatz Texture
          Dataset original no incluye etiquetas categóricas. Sin embargo, existen versiones modificadas
          y extensiones del dataset que sí incluyen etiquetas o puedes asignar etiquetas basadas en la
          tabla de referencia.
          \begin{itemize}
              \item Mostrar imagen e histograma de una imagen de cada clase.
          \end{itemize}
    \item \textbf{Implementación de LBP.} El algoritmo LBP sigue estos pasos:
          \begin{enumerate}
              \item Para cada píxel central, comparar su intensidad con 8 vecinos.
              \item Generar un patrón binario (1 si el vecino > central, 0 en otro caso).
              \item Convertir el binario a decimal y reemplazar el píxel central.
          \end{enumerate}
    \item \textbf{Extracción de Características:} Convertir cada imagen LBP en un histograma de
          256 bins (para valores 0-255).
          \begin{enumerate}
              \item Mostrar el histograma normalizado de una imagen de cada clase.
              \item Comparar histogramas por clase (ej. madera vs. tela vs. piedra).
          \end{enumerate}
    \item \textbf{Clasificación y Evaluación}
          \begin{enumerate}
              \item Divide los datos en entrenamiento (70\%) y prueba (30\%).
              \item Usa un modelo simple como vecinos más cercanos (k-NN) o Máquina de Soporte Vectorial
                    (SVM) sobre los vectores de características y entrena el modelo.
              \item Calcula la exactitud en el conjunto de prueba.
              \item Visualiza las imágenes LBP para entender qué patrones se están detectando.
              \item Mostrar matriz de confusión.
          \end{enumerate}
    \item \textbf{Visualización de Resultados:}
          \begin{enumerate}
              \item Descripción de los pasos realizados.
              \item Resultados obtenidos (gráficos y análisis).
          \end{enumerate}
    \item Repite los pasos 3, 4, 5 y 6 aplicando el algoritmo LBP con configuraciones como
          $P=16, R=2$ y $P=24, R=3$ y compara resultados.
\end{enumerate}